\section{関連技術}
\label{sec_related}
本章では，本研究で比較した手法の基礎となる技術を整理する．

\subsection{画像とテキストの埋め込みモデル}
Contrastive Language-Image Pre-training (CLIP)\cite{ref_clip} は，画像と文章を同じベクトル空間に写す2つのエンコーダを，大量の画像と説明文の組から対照学習したモデルである．
以下では，エンコーダが出力するベクトルを埋め込みと呼ぶ．
SigLIP\cite{ref_siglip} は，CLIP の対照学習の損失をシグモイド関数による2値の損失に置き換えたモデルであり，SigLIP 2\cite{ref_siglip2} は多言語の学習データや位置の情報を扱う学習を加えた後継のモデルである．
本研究では，CLIP の ViT-B/32 のモデル（以下，CLIP-B/32），入力が224画素の SigLIP 2 の Base のモデル（以下，SigLIP2-B），入力が384画素の SigLIP 2 の so400m のモデル（以下，SigLIP2-so400m）を用いた．

埋め込みモデルによる分類には2通りの使い方がある．
1つ目はゼロショット分類であり，各クラスを表す文章の埋め込みと画像の埋め込みのコサイン類似度が最大のクラスを選ぶ．
学習データは不要であるが，分類の精度は文章の選び方に依存する．
2つ目は，画像の埋め込みを特徴量とし，ラベル付きのデータでロジスティック回帰などの分類器を学習する方法である．
特に線形の分類器を用いる場合を線形プローブと呼ぶ．
いずれも1枚あたりの計算量は画像エンコーダの1回の順伝播のみであり，GPU で毎秒数百から千枚程度を処理できる．

\subsection{視覚言語モデル}
VLM は，画像と指示文を入力として文章を生成するモデルである．
クラスの定義を自然言語で与えるだけで分類でき，学習データを必要としない．
本研究では，Qwen3-VL\cite{ref_qwen3vl} の Instruct のモデルを主に用い，比較として InternVL3.5\cite{ref_internvl35} と MiniCPM-V 4.5\cite{ref_minicpmv} を用いた．
推論には，大規模言語モデルの推論を高速化するライブラリ vLLM\cite{ref_vllm} を用いた．
vLLM の構造化出力の機能を用いると，出力を JSON スキーマや正規表現に一致する文字列に制約できる．
本研究では全ての VLM の出力をラベルの候補に制約したため，出力を解釈できない事例は生じなかった．
VLM は精度が高い一方で，1枚あたりの計算量が埋め込みモデルより1桁以上大きく，Qwen3-VL-8B で毎秒25枚程度であった．

\subsection{知識蒸留}
知識蒸留\cite{ref_kd}は，大きく高精度なモデル（教師）の出力を目標として，小さく高速なモデル（生徒）を学習する方法である．
本研究では，教師である VLM が一部の画像に付けたラベルを正解として，埋め込みに基づく分類器や画像エンコーダを学習する．
教師が全件を処理する代わりに，教師は数万枚にラベルを付けるのみであり，全件の処理は高速な生徒が行う．
また，埋め込みを一度計算して保存しておけば，ラベルの定義を変えた場合も，教師ラベルを付け直して小さな分類器を再学習するだけで全件のラベルを更新できる．

\subsection{画像エンコーダの微調整}
埋め込みを固定して分類器のみを学習する代わりに，画像エンコーダ全体を教師ラベルで追加学習することもできる．
以下では，これを微調整と呼ぶ．
微調整は，埋め込みに基づく分類器より多くの計算を学習時に要するが，小さな画像エンコーダでも大きなモデルの埋め込みに近い精度を得られる場合がある．
本研究では，SigLIP2-B と ConvNeXt-Tiny\cite{ref_convnext} の微調整を比較した．
